\documentclass[10pt]{article}
\usepackage{template}
\usepackage{amsmath,amssymb,tikz,graphicx,xurl,microtype}
\usetikzlibrary{arrows.meta,positioning,calc,decorations.pathreplacing}
\setlength{\emergencystretch}{3em}
\raggedbottom
\SetKeys{date=2026-11-03,vol=021,title={每日有趣新闻｜选票、证明与证据},author={ChatGPT 编辑部},note={\small 采编截止：北京时间11月3日上午。主要核查11月1—2日新事实；数学取11月2日arXiv新列表，并逐一注明原始提交日期。本文区分媒体报道、官方声明、作者主张及本刊分析。}}
\renewcommand{\newsitem}[3]{\par\medskip{\large\bfseries\color{newsblue}#1\par}{\footnotesize\color{newsgray}#2\par}\smallskip\color{black}#3\par\medskip{\color{newsgray}\rule{\linewidth}{0.25pt}}\par}
\newcommand{\newslead}[3]{\par\medskip{\LARGE\bfseries\color{newsblue}#1\par}\vspace{1mm}{\small\color{newsgray}#2\par}\smallskip\color{black}#3\par\medskip{\color{newsgray}\rule{\linewidth}{0.45pt}}\par}
\newcommand{\newssub}[1]{\par\smallskip\noindent{\bfseries\color{newsblue}#1}\quad}
\begin{document}
\begin{documents}
\raggedcolumns

\newslead{美国中期选举前夜：席位预测不是选举结果，法律程序也不是计票软件}{国际政治深度｜11月2日路透社选举指南、路透／益普索民调；11月3日投票}{
美国将在11月3日举行中期选举，全部众议院席位、部分参议院席位以及多州州长职位将接受选民检验。总统本人不在这张选票上，但国会两院的归属决定了未来两年行政当局能否较顺畅地推进预算、立法与任命。路透社11月2日的选举指南称，共和党在参议院拥有53比47的优势；民主党需要净增至少四席才有望翻转参议院。众议院当时的席位分布为218比214，仍有空缺，因此同样不能把两党席位简单相加当成435席全员在位。\textbf{所有这些都是投票前的制度与席位背景，不是11月3日已经出现的结果。}\cite{ReutersElection,ReutersSenate}

\newssub{为什么选举会出现「赢得更多票，却没有掌握国会」？}众议院按选区分别选出，参议院每州两席，人口大小不改变州的参议院代表数。全国政党支持率不是将选票倒进一个桶里按比例分配席位。假设一个政党在许多选区以很大优势胜出，却在更多选区以微弱差距落败，其全国总票数可能领先而席位数落后。美国的选区划分、各州人口和投票率会进一步改变这种映射。民调中所谓「通用国会选票」只是在问选民更倾向哪个党，不等于对每一个竞争选区的精确预测。\cite{ReutersElection,ReutersPoll}

\begin{center}\begin{tikzpicture}[x=0.085cm,y=0.35cm]
\node[anchor=west,font=\scriptsize] at (0,2.7) {参议院：投票前席位（共100席）};
\fill[newsblue!75] (0,1.4) rectangle (47,2.2);
\fill[red!70!black] (47,1.4) rectangle (100,2.2);
\node[white,font=\scriptsize] at (23.5,1.8) {民主党47};
\node[white,font=\scriptsize] at (73.5,1.8) {共和党53};
\draw[very thick,dashed] (51,1.2)--(51,2.5);
\node[font=\tiny,anchor=west] at (0,.4) {虚线：取得多数席所需的51席位置；不是民调预测。};
\end{tikzpicture}\par{\tiny\color{newsgray}图1｜按路透社11月2日投票前统计重绘；参议院席位分布，非选举结果。}\end{center}\color{black}

\newssub{一张民调表能告诉我们什么？}路透社／益普索11月2日公布的调查将特朗普施政支持率列为32\%，在通用国会选票问题上民主党与共和党分别获得46\%和38\%的支持。这说明受访者总体倾向，但不能据此把八个百分点直接换算成八个百分点的席位差：首先抽样存在误差；其次未决定选民、不同州选民参与率及实际候选人都影响结果；最后，选区票数与席位之间本就不存在固定的线性关系。与其问「民调准不准」，更值得问它测量的对象、抽样时间和所能预测的变量究竟是什么。\cite{ReutersPoll}

\newssub{投票、计票、认证，是三道不同的门。}路透社报道，司法部计划在选举日部署约一千名观察人员，比上届中期选举更多。司法部称此举有助监督依法投票，一些地方选举官员则担忧联邦人员的介入可能形成压力。对于这类争议，不能仅凭部署人数判定选举被操纵，也不能因为官方称为「监督」就排除程序合法性的审查。实际计票由各州和地方依法律进行，媒体根据已统计票数与剩余选票作出「胜负判定」，并不等于法定认证；诉讼、重新计票、认证和就职各有独立时间表。\cite{ReutersElection,ReutersCalls}

\newssub{什么情况下新闻媒体会提前宣布胜者？}路透社11月1日解释，美国新闻机构通常使用已公布计票、提前投票和历史数据估计尚未统计的选票构成。只有在剩余选票不足以逆转差距时，才会宣布候选人获胜。因而一场选举可能在几小时内得到高度可信的媒体判定，却仍要等待官方确认；反过来，部分州可能数日无法确定胜负，这也不自动证明出现舞弊。\cite{ReutersCalls}

\newssub{政治含义并不止于「红」或「蓝」。}路透社列举阿拉斯加、堪萨斯、缅因、北卡罗来纳、俄亥俄、得克萨斯和密歇根等关键竞争州。民主党如果取得一院控制权，可能加强对行政部门政策、预算与战争决策的调查；共和党如果维持两院优势，行政议程则可能获得更大的立法空间。两党内部也会从结果中争论2028年应当转向温和派还是动员更激进的选民。此处都是基于制度与竞选格局的\textbf{条件推断}，不是本刊对哪一党应当获胜的倡议。\cite{ReutersElection,ReutersSenate}

\newssub{编辑观察：最该警惕的是时间错位。}11月3日上午（北京时间）距离美国本土大多数投票站开放仍有时差；此时传播的「某党已赢得众议院」若没有明确的官方或可信媒体出处，尤其值得警惕。接下来应跟踪各州实际投票人数、提前票与当日票的统计差异、诉讼进展及认证结果。对任何阵营而言，程序的可验证性比选举夜的即时口号更重要。\cite{ReutersElection,ReutersPoll}
}

\newslead{一次求根能问函数几个问题？Kung--Traub猜想的新证明与Lean核验}{数学与形式化深度｜Colbrook；arXiv:2610.40660；11月2日数学新列表，原稿9月20日提交}{
数学家想用尽可能少的函数计算找到方程$f(x)=0$的解。牛顿法每步求一次函数值$f(x)$与一次导数$f'(x)$，再作$x_{k+1}=x_k-f(x_k)/f'(x_k)$。在根附近且根为单根时，它通常具有二阶收敛：若误差为$e_k=|x_k-\alpha|$，则下一步误差近似按$e_{k+1}\le C e_k^2$缩小。问题随之出现：如果每轮最多向函数询问$n$个实数值（包括任意阶导数），又允许算法自适应决定下一问的位置，能否构造出对所有足够光滑的函数都比$2^{n-1}$阶更快的求根方法？\cite{KungPaper,KungCode}

\newssub{先定义「阶」，再讨论「最优」。}设$\alpha$是$f$的一个单根，即$f(\alpha)=0$且$f'(\alpha)\ne0$。算法$A$在根附近具有局部$p$阶收敛，意思是存在$C,\delta>0$，对每个满足$0<|x-\alpha|<\delta$的起点，都有$|A(f,x)-\alpha|\le C|x-\alpha|^p$。这是一个\textbf{对所有充分接近的起点同时成立的界}，不能仅用某条数值迭代轨迹上的拟合斜率替代。更关键的是「普遍最优」：这里讨论的是算法面对所有具有单根的实解析函数，而不是它在某个特别好算的多项式上表现多快。\cite{KungPaper,KungCode}

\begin{center}\begin{tikzpicture}[x=1cm,y=0.5cm]
\draw[->] (0,0)--(6.9,0) node[right,font=\tiny] {每轮观测次数 $n$};
\draw[->] (0,0)--(0,4.7) node[above,font=\tiny] {阶数};
\foreach \n/\v in {1/1,2/2,3/4,4/8}{\fill[newsblue!80] ({1.2*\n-0.45},0) rectangle ({1.2*\n+0.05},{0.45*\v});\node[font=\tiny] at ({1.2*\n-0.2},{0.45*\v+0.24}) {$\v$};\node[font=\tiny] at ({1.2*\n-0.2},-0.24) {$\n$};}
\end{tikzpicture}\par{\tiny\color{newsgray}图2｜根据 $2^{n-1}$ 绘制的理论上界示意，纵轴按线性比例；并非实验性能数据。}\end{center}\color{black}

\newssub{猜想到底声称什么？}Kung和Traub在1974年的工作提出，若一轮只有$n$次函数或导数的标量观测、方法不保留跨轮记忆，那么可达到的普遍局部阶数至多为$2^{n-1}$。例如$n=2$给出二阶，$n=3$给出四阶。Matthew J. Colbrook的新论文主张证明了这个上界，并指出反插值的Hermite方法可以达到它，因此上界不是随便取的保守数字。对方法的查询位置与导数阶数可以自适应，计算规则不要求连续、可测或可计算；限制主要来自每轮可取得的独立观测数量。\cite{KungPaper,KungCode,KungTraub1974}

\newssub{难点在「存在一个足够坏的函数」。}为了否定所有可能算法都能超过某个阶数，不能只展示牛顿法或割线法的弱点。作者利用多项式Wronskian的局部化性质，在任意固定尺度上构造让观测无法充分区分的函数，再以对角化方式拼成一个整函数。整函数在整个复平面解析，限制到实轴后仍足够规则；仓库说明中的主要反例甚至有唯一实根，并且在实轴上满足$3/4\le f'(t)\le5/4$。也就是说，反例并非依赖一个奇形怪状、不可微的病态函数。\cite{KungPaper,KungCode}

\newssub{Lean能证明什么，不能证明什么？}作者提供的公开仓库固定Lean 4.33.1与Mathlib版本，包含主要实数结论、分组观测推论、复数版本及达到上界的构造。README记载发布快照通过了干净构建、1202项声明的公理审计、1058个具名定理的比较检查，并另列12项主要结果的独立检查记录。这些数字是\textbf{仓库维护者报告的机器检查结果}，本刊未在独立环境重跑全部检查。形式化证明若确实由Lean内核接受，能排除许多推理层面的疏漏；但形式化定理的定义是否忠实对应论文里的自然语言问题，仍需数学家审查。\cite{KungCode,KungPaper}

\newssub{验证的边界。}原论文在arXiv的版本标记为9月20日提交；它在11月2日的数学新列表出现，不应被写成11月2日首次发现。仓库注明ChatGPT 6参与Lean形式化，作者审阅并承担数学责任。这一披露不是结果自动正确的证据，也不是可以忽略形式化记录的理由。当前能确认的是：有具体证明文本、可查的形式化仓库和维护者的构建审计说明；尚无本刊独立复核或已获期刊同行评议的依据。对喜欢形式化证明的人，这比「AI又解决了一个猜想」的标题更值得细读。\cite{KungPaper,KungCode}
}

\newsitem{多哈会谈结束却没有突破：海峡「两天方案」到底改变了什么？}{国际外交重点｜11月2日路透社、美联社；与昨日报道相比有新进展}{
11月2日，伊朗外交部表示，与美国之间通过卡塔尔等调解方传递的信息仍在继续，但此前几天在多哈的接触没有取得突破。路透社报道，伊朗代表已于周日离开多哈，调解人随后赴德黑兰讨论草案。美联社则记录了美国副总统万斯对达成协议的谨慎乐观，以及以色列领导层敦促美国在中期选举后重新施压的消息。双方的口径不同，却都不能证明停火已经签署。\cite{ReutersIran2,APIran2}

\newssub{七步方案被压缩，不代表条件被取消。}伊朗发言人称，九月提出的七步方案可被压缩到两天内执行，但前提仍包括美国解除港口封锁、伊朗恢复霍尔木兹海峡通航、油品销售限制及制裁问题。伊朗强调，只有先核验相关条件落实，才可能重新讨论核计划；美国方面则要求伊朗作出可验证的不重建核计划承诺。这是谈判顺序的冲突，不是仅凭缩短日程就能消除的技术问题。\cite{ReutersIran2,APIran2}

\begin{center}\begin{tikzpicture}[>=Stealth, node distance=4mm, every node/.style={font=\scriptsize}]
\node[draw,rounded corners,fill=newsblue!9,minimum width=2.9cm,align=center] (a) {美国港口封锁\\与制裁安排};
\node[draw,rounded corners,fill=newsblue!9,minimum width=2.9cm,align=center,right=8mm of a] (b) {伊朗海峡通航\\与核计划条件};
\draw[<->,very thick,newsblue] (a)--node[above,font=\tiny] {谈判次序争议} (b);
\node[below=5mm of a,align=center,font=\tiny] {解除／核验\\尚未完成};
\node[below=5mm of b,align=center,font=\tiny] {恢复／核验\\尚未完成};
\end{tikzpicture}\par{\tiny\color{newsgray}图3｜据路透社与AP 11月2日报道归纳的条件关系；非正式协议流程图。}\end{center}\color{black}

\newssub{为什么石油价格下跌不能当作和平证据？}路透社称，11月2日油价因外交希望一度下跌近4\%。期货价格反映市场参与者对未来供需的预期，并不是海峡实际开放的船舶计数器。与此同时，一艘卡塔尔液化天然气运输船此前遭不明物体击中，运营商未认定攻击者。任何将这次袭击归咎于某国的断言都需要额外证据。接下来应看正式协议文本、港口通行指令、独立航运数据与核查机制，而不是只盯着外交官的形容词。\cite{ReutersIran2,APIran2}
}

\newsitem{中美农业采购的算术：72亿美元与170亿美元，能直接判定违约吗？}{中国经济与贸易重点｜路透社11月2日海关数据分析；白宫五月承诺口径}{
路透社根据海关数据计算，中国2026年1—9月购买了约72亿美元的美国非大豆农产品，是2019年以来同期最低水平；而美国白宫在5月称，中国承诺自2026年起每年采购约170亿美元此类商品，另有单独的大豆采购安排。乍看之下，72与170的差距巨大，但不能把前九个月的实绩直接同全年目标相减，就宣布中国已经违约。白宫尚未明确承诺计时究竟从五月公告算起，还是从有关关税措施正式落实后开始；该问题尚未得到中美双方的清晰解释。\cite{ReutersFarm}

\begin{center}\begin{tikzpicture}[x=0.30cm,y=0.52cm]
\node[anchor=west,font=\scriptsize] at (0,2.4) {美国非大豆农产品对华采购（亿美元）};
\fill[newsblue!65] (0,1.2) rectangle (7.2,1.8);
\fill[newsblue!12,draw=newsblue] (0,.2) rectangle (17,.8);
\node[anchor=west,font=\tiny] at (7.4,1.5) {72（1—9月实绩）};
\node[anchor=west,font=\tiny] at (17.2,.5) {170（白宫称年度承诺）};
\end{tikzpicture}\par{\tiny\color{newsgray}图4｜路透社11月2日海关数据分析与白宫承诺口径；两项时间范围不同，不构成履约率。}\end{center}\color{black}

\newssub{需求不是签约时按一下按钮就能创造的。}玉米、小麦、高粱、肉类都受到国内库存、国际报价、物流成本和检疫许可约束。路透社举例称，中国2021年从美国进口玉米价值约56亿美元，但国内供应增加后，美国玉米竞争力减弱；巴西高粱也在争夺中国市场。猪肉、鸡肉供应充裕，美国牛肉则受价格和部分出口资格限制。即使两国政府希望增加采购，企业仍需考虑利润与风险。\cite{ReutersFarm}

\newssub{政治承诺与贸易统计的边界。}有分析人士认为若要达到目标，可能需要国有买家承担更多采购；这属于受访者的政策推断，并非中国政府已经下达新采购命令。报道还指出美国农民面临库存与成本压力，中期选举使农业贸易问题具有明显政治含义。要判断履约情况，应继续核对完整协议文本、起算日期、商品范围、年度海关数据及双方对差异的正式解释，而不能把媒体分析的72亿美元说成中国官方承认的欠款。\cite{ReutersFarm}
}

\newsitem{哈勃发现矮星系「缺氧」：元素不是凭空消失，而是可能被恒星风吹走}{天文学重点｜NASA 11月2日发布；《天体物理学杂志》当日论文}{
NASA 11月2日介绍一项名为GLOW（Galaxies Losing Oxygen via Winds）的研究。团队使用哈勃太空望远镜的档案光学观测，结合斯皮策、甚大阵列及澳大利亚望远镜紧凑阵列的数据，研究37个附近矮星系，估计它们在恒星演化中制造了多少氧，又保留了多少。结果显示，许多矮星系所保留的氧比模型预期少得多。\textbf{「缺氧」指的是星系化学演化的物质收支，而不是发现星系中的氧原子违反守恒定律。}\cite{NASAOxygen,OxygenPaper}

\newssub{为什么选氧而不是铁？}氧是恒星核合成的重要产物，在天文光谱中有便于识别的谱线，因而可作为追踪恒星形成与元素丰度的指标。大质量恒星、超新星和星系风能够把富含重元素的气体从引力较弱的矮星系吹向周围空间。团队由恒星群体的颜色、亮度重建恒星形成历史，再与气体、尘埃中的氧含量比较，发现实际保留比例偏低。\cite{NASAOxygen,OxygenPaper}

\begin{center}\begin{tikzpicture}[>=Stealth,node distance=5mm,every node/.style={font=\scriptsize}]
\node[draw,rounded corners,fill=orange!13,align=center] (s) {恒星核合成\\产生氧};
\node[draw,rounded corners,fill=newsblue!10,align=center,right=9mm of s] (g) {矮星系气体\\与尘埃};
\node[draw,rounded corners,fill=green!10,align=center,right=9mm of g] (o) {星系周围气体晕\\或星系际空间};
\draw[->,thick] (s)--(g);
\draw[->,thick] (g)--node[above,font=\tiny] {星系风／超新星} (o);
\end{tikzpicture}\par{\tiny\color{newsgray}图5｜根据NASA GLOW项目11月2日介绍绘制的元素迁移示意；未定量测出氧的最终去向。}\end{center}\color{black}

\newssub{研究还没回答「氧去了哪里」。}氧可能停留在星系周围稀薄的气体晕，也可能已经进入星系际介质。NASA说明目前尚不清楚两者各占多少，因此不能把「模型预测与观测有差异」直接写成「已找到所有流失氧的藏身之处」。样本只有37个、均为足够近而能分辨恒星的矮星系；对更远更小星系的推广仍需新观测。论文已发表于同行评议期刊，但本刊未独立重做丰度计算。NASA同时公开NGC 4449的哈勃原始图像及署名，本期因无法取得可归档图像文件，改以明确标注的元素流向示意图呈现，不将示意图伪装为天文照片。\cite{NASAOxygen,OxygenPaper}
}

\newsitem{安全团队早在2024年预警：日本GMO外汇交易服务遭入侵后的治理问题}{网络安全与企业治理重点｜11月2日东京记者会，路透社日文报道}{
11月2日，日本GMO Internet Group就外汇保证金（FX）自动交易服务遭未经授权访问一事举行记者会。集团负责人熊谷正寿致歉，并承认集团旗下安全专业子公司早在2024年就提出过问题，但这些警告没有转化为运营措施；据其说法，警告当时也没有传达到集团最高负责人。这一新信息把事件从「技术人员没打补丁」推进到组织如何处理内部预警的问题。\cite{ReutersGMO}

\newssub{预警为何可能在组织中失效？}大型集团有多层子公司、产品线和安全团队。漏洞报告可能进入工单系统，却没有对应的责任人、期限、复验流程与向上升级条件；即使技术风险被识别，预算与业务连续性压力也可能推迟整改。需要强调的是，路透社报道的是管理层对过去警告的承认，并未公开全部入侵取证材料。本刊不能据此断言某个具体CVE就是攻击入口，也不能猜测用户账户、资金损失或数据外泄规模。\cite{ReutersGMO}

\newssub{可检验的整改标准。}企业若宣称将建设「网络防御」体系，应进一步公开问题分类、预警到整改的闭环时长、独立复验结果，以及受影响客户的补救方式。技术安全不是每年写一次报告，而是把发现、处置、验证和问责串成有记录的流程。这也是同一天媒体报道中最值得软件工程师关注的部分。\cite{ReutersGMO}
}

\newsitem{奥列什基六日停火结束：近八百人撤离之后，双方又在争论谁利用了空档}{俄乌战争与人道救援重点｜11月2日路透社；乌方、俄方及美方说法分别标注}{
乌克兰南部奥列什基的局部人道停火于11月1日夜结束。美国白宫称，停火期间有近八百名居民从俄方控制区撤离，双方此前均对人道行动的完成表示肯定。11月2日，乌克兰总参谋部却指责俄军趁停火轮换部队、加固防御、改善后勤；克里姆林宫发言人佩斯科夫表示没有掌握具体情况，但称停火期间部队调整属于通常做法，并反问乌方是否也采取类似行动。\textbf{双方都承认存在停火与撤离行动，并不意味着对军事行为是否违反协议达成一致。}\cite{ReutersOleshky}

\newssub{救出多少人，为什么还不能算出行动成功率？}「近八百」是美方给出的撤离人数，并非战前全镇人口，也不是原先登记愿意撤离的人数。俄方人权专员表示镇内仍有不到750人，并称将安排后续定期援助。由于不同来源可能采用不同登记时间、居民范围和迁移路线，不能把两组数字相减就断定有多少人失踪。真正应当跟踪的是每名撤离者是否自愿选择目的地、医疗救助是否持续，以及留守者是否仍能获得基本物资。\cite{ReutersOleshky}

\newssub{同一天，另一处救援点遭袭。}路透社援引乌克兰总统和紧急救援部门称，扎波罗热州科梅舒瓦哈一所用作避难与每周人道物资发放点的学校遭俄军航空炸弹袭击，造成至少26人死亡，其中包括救援人员与警察；另有失踪者。上述死亡人数、攻击方和两轮袭击的细节来自乌方官员，报道时应注明其归属；本刊尚无独立现场弹道鉴定或俄方对该具体事件的完整回应。平民保护与调查责任必须严肃对待，但严肃不等于跳过证据核查。\cite{ReutersZaporizhzhia}

\newssub{局部人道停火不是全面停战。}停火的地理范围、期限、可允许的行动以及监督机构，都决定了「违约」的法律和事实判断。即便撤离确实完成，其他地区仍可能遭袭；即便有军事重新部署，也需核对双方当初约定是否明确禁止。路透社另引述美国高级官员称，俄乌代表计划于11月10—11日在阿布扎比讨论降级与战后安全安排。这是预定会谈，而非已经达成的新和平协议。\cite{ReutersOleshky,ReutersZaporizhzhia}
}

\newsitem{中国外汇存款准备金率拟升至6\%：央行为什么在人民币走强时回收美元流动性？}{中国货币政策重点｜11月3日凌晨路透社与中国人民银行公告口径}{
中国人民银行11月3日宣布，自11月15日起将金融机构外汇存款准备金率由4\%提高至6\%，上调两个百分点。路透社根据外汇存款数据估算，新增准备金需求可能从金融体系回收约240亿美元的可用外汇流动性。这是估算的流动性影响，不是央行直接在外汇市场抛售240亿美元，也不是人民币法定存款准备金率同步变化。\cite{ReutersFX}

\begin{center}\begin{tikzpicture}[x=.30cm,y=.65cm]
\node[anchor=west,font=\scriptsize] at (0,2.2) {每100美元外汇存款对应准备金（美元）};
\fill[newsblue!35] (0,1.15) rectangle (4,1.75);
\fill[newsblue!75] (0,.2) rectangle (6,.8);
\node[anchor=west,font=\tiny] at (4.3,1.45) {调整前4美元};
\node[anchor=west,font=\tiny] at (6.3,.5) {11月15日起6美元};
\end{tikzpicture}\par{\tiny\color{newsgray}图6｜据11月3日央行公告比例绘制的等规模示意；不是市场美元存量统计。}\end{center}\color{black}

\newssub{从资产负债表看机制。}银行吸收客户美元存款后，并非可以把每一美元都用于贷款或交易；一部分须按规定缴存外汇准备金。准备金率从4\%升至6\%，在存款规模不变的理想化情况下，每100美元外汇存款需多留2美元作为准备金，市场上可自由配置的美元减少。其他条件相同时，美元流动性收紧可能缓和人民币升值压力，但汇率还受到利差、贸易顺差、资本流动和预期影响，不能把政策效果写成精确的汇率预测。\cite{ReutersFX}

\newssub{为何是现在？}路透社称，人民币今年以来兑美元累计升值逾5\%，出口与贸易顺差、企业结汇是市场背景。2023年9月央行曾下调外汇存款准备金率以应对当时的人民币贬值压力；本次方向相反。官方表述是加强金融机构外汇流动性管理，分析师则认为意在降低市场对单边升值的预期。前者是政策公告，后者是对政策动机的解释，不能混成一条「央行承认要打压人民币」的引语。\cite{ReutersFX}
}

\newsitem{英国网安中心警告AI增强型攻击：最值得担心的不是「自动黑客人格」}{网络安全与AI重点｜11月3日英国NCSC年度报告，路透社}{
英国国家网络安全中心（NCSC）在最新年度报告中表示，未来两年针对关键基础设施出现AI增强的攻击具有高度现实风险；其判断包括到2028年前可能发生使用AI增强恶意软件的勒索攻击。NCSC负责人Richard Horne强调，许多组织和供应链长期存在的技术欠账，正使攻击者能够借助AI更快发现并利用旧漏洞。\textbf{这是一份面向未来的风险评估，不是英国已经确认发生了某次完全自主的AI攻击。}\cite{ReutersNCSC}

\newssub{为什么旧漏洞比新名词更重要？}传统攻击链需要资产探测、漏洞匹配、凭据窃取、横向移动、勒索与数据外传等环节。AI可以加速脚本编写、钓鱼文本变体生成与海量告警筛选，但是否能在真实环境稳定完成端到端攻击，取决于目标防护、访问权限和执行反馈。将AI描述为凭空创造新的网络攻击物理规律并不准确；更可检验的危险是相同攻击成本下降、速度加快，使原本可以在数周内修复的漏洞更快被利用。\cite{ReutersNCSC}

\newssub{归因必须保持距离。}NCSC报告对中国、俄罗斯、伊朗等国家行为体的网络能力作出安全评估，其中对中国的表述尤其严厉。这是英国政府机构基于本国国家安全职责形成的判断，并非本刊独立完成的归因，也不应把某国公民、企业与国家情报活动混同。对普通运维团队，最有价值的建议仍是建立资产清单、限制管理面暴露、分离凭据、及时更新、备份演练和事件响应记录。AI安全不是替代这些基础措施的魔法工具。\cite{ReutersNCSC}
}

\newsitem{数学拾趣｜阿拉巴马悖论：议会多加一席，一个州却可能少一席}{数学史与组合几何｜David McCune；arXiv:2610.40682；10月26日提交、11月2日新列表}{
席位按人口比例分配，听起来很自然；可席位必须是整数。\textbf{阿拉巴马悖论}说的是：如果按某种「最大余数法」分配席位，在各地区人口不变时，增加议会总席位，反而可能使某个地区少拿一席。例子：三个地区人口比例为$5:3:1$。总共4席时，各自理论份额为$20/9,12/9,4/9$，先取整数得$2,1,0$，剩下一席给余数最大的第三地区，结果$2,1,1$；总共5席时理论份额为$25/9,15/9,5/9$，先取整数得$2,1,0$，剩两席给余数较大的第二和第一地区，结果$3,2,0$。第三地区在总席位增加时竟从1席变为0席。\cite{ApportionPaper}

\newssub{为什么不矛盾？}最大余数法并不是「每个地区席位随总席位单调增加」的算法，而是每次重新计算所有小数余数的排名。总席位增加会改变每个地区理论份额的小数部分，排名可能突然交换。David McCune的新预印本重新审视1992年一篇几何分析，声称修正其中的概率公式，并讨论相关几何后果。它是一篇准备投稿的解释性论文，尚无本刊独立核验或已完成同行评议的证据。这里的数字例子是本刊独立演算的教学示意，不是作者原文数据。\cite{ApportionPaper}
}

\newsitem{简讯｜汇丰计划在香港设AI研究所}{金融科技｜11月2日香港金融科技周，路透社}{
汇丰行政总裁在11月2日香港金融科技周宣布，计划于2027年第一季度在香港设立人工智能与新兴技术研究所，以香港科技大学为主要学术伙伴。此前汇丰已在新加坡设立类似机构。公司尚未公布人员规模；「计划开设」不等于研究所已经投入运行，也不代表任何具体模型成果已经产生。\cite{ReutersHSBC}
}

\newsitem{简讯｜旅行者1号离「一光日」越来越近}{航天工程｜NASA10月29日通报，11月18日预计达到距离里程碑}{
NASA披露，旅行者1号在9月完成三个耗电设备的同步替换，10月18—19日两次计划通信未收到信号，20日自行恢复。团队确认推进器组发生切换，但仍在调查故障原因。NASA预计它11月18日到达距地球一光日的位置：单程无线电传播约24小时，往返约48小时。它不会在那一刻发生物理跃迁；这只是距离标志。每年约4瓦的电力损失仍是任务长期限制。此消息的原始发布日期为10月29日，本刊将其作为航天工程补充，而非冒充11月2日的新故障。\cite{NASAVoyager}
}

\par\medskip{\large\bfseries\color{newsblue}采编、图片与核查说明}\par
{\scriptsize 本期核查11月1—2日国际时政、战争与和平、经济政策、社会、科技与AI、Rust/C++、Linux、物理、化学、航天、文化，以及11月2日arXiv数学新列表（数论、代数、几何、组合、逻辑、概率与Lean/mathlib）。选题按可靠增量取舍，不强制每类刊登。共六张原创TikZ图，分别为席位示意、求根阶数、谈判条件、农产品采购口径、矮星系元素流向和外汇准备金率；全部图源随正文归档。NASA公布了NGC 4449可核查的原始哈勃照片，图像署名为NASA/ESA/Hubble Heritage/Alessandra Aloisi，但本次无法取得原图文件，因此不冒充已归档照片；路透社/AP现场照片需要另行授权，未擅自复制。数学预印本均区分作者主张、机器检查及同行评议状态。}\par
\begingroup\fontsize{6}{7.2}\selectfont
\renewcommand{\refname}{参考文献}
\bibliographystyle{plainurl}\bibliography{refs}
\endgroup
\end{documents}
\end{document}